\documentclass[12pt,a4paper]{article}

% ─── Pacotes ────────────────────────────────────────────────────────────────
\usepackage[utf8]{inputenc}
\usepackage[T1]{fontenc}
\usepackage[brazil]{babel}
\usepackage{geometry}
\usepackage{graphicx}
\usepackage{booktabs}
\usepackage{array}
\usepackage{tabularx}
\usepackage{amsmath}
\usepackage{amssymb}
\usepackage{hyperref}
\usepackage{xcolor}
\usepackage{setspace}
\usepackage{parskip}
\usepackage{enumitem}
\usepackage{float}
\usepackage{caption}
\usepackage{fancyhdr}
\usepackage{titlesec}

% ─── Geometria da página ────────────────────────────────────────────────────
\geometry{top=2.5cm, bottom=2.5cm, left=3cm, right=2.5cm}

% ─── Cabeçalho e rodapé ─────────────────────────────────────────────────────
\pagestyle{fancy}
\fancyhf{}
\rhead{\small POC II --- Proposta de Projeto}
\lhead{\small UFMG / DCC}
\cfoot{\thepage}

% ─── Formatação de seções ───────────────────────────────────────────────────
\titleformat{\section}{\large\bfseries}{\thesection}{1em}{}[\titlerule]
\titleformat{\subsection}{\normalsize\bfseries}{\thesubsection}{1em}{}
\titleformat{\subsubsection}{\normalsize\itshape}{\thesubsubsection}{1em}{}

% ─── Espaçamento ────────────────────────────────────────────────────────────
\setlength{\parskip}{6pt}
\onehalfspacing

% ─── Hyperlinks ─────────────────────────────────────────────────────────────
\hypersetup{colorlinks=true, linkcolor=blue!70!black, citecolor=blue!70!black, urlcolor=blue!70!black}

% ════════════════════════════════════════════════════════════════════════════
\begin{document}

% ─── Capa ───────────────────────────────────────────────────────────────────
\begin{titlepage}
  \centering
  \vspace*{1cm}
  {\large\textbf{Universidade Federal de Minas Gerais}}\\[0.3cm]
  {\large Departamento de Ciência da Computação}\\[0.2cm]
  {\normalsize Belo Horizonte, MG, Brasil}
  \vspace{3cm}

  {\LARGE\bfseries Proposta de Projeto Orientado a Computação~II}

  \vspace{1.5cm}

  {\large\bfseries
    Handover Inteligente em Redes LEO com\\
    Reward Shaping Multiobjetivo, Memória\\
    Temporal (LSTM) e PPO
  }

  \vspace{2cm}

  \begin{tabular}{ll}
    \textbf{Aluno:}      & Sacramento \\[4pt]
    \textbf{Orientador:} & --- \\[4pt]
    \textbf{Semestre:}   & 2026/2 \\
  \end{tabular}

  \vfill
  {\normalsize Belo Horizonte, \today}
\end{titlepage}

% ─── Sumário ────────────────────────────────────────────────────────────────
\tableofcontents
\newpage

% ════════════════════════════════════════════════════════════════════════════
\section{Introdução}

A integração de redes via satélite é um pilar fundamental da sexta geração de
conectividade (6G), com o objetivo de fornecer cobertura global e ininterrupta.
Constelações de satélites \textit{Low Earth Orbit} (LEO), operando a
aproximadamente 550\,km de altitude, oferecem baixa latência, mas impõem o
desafio de handovers frequentes --- a cada 2 a 5 minutos --- devido à alta
velocidade orbital ($\approx 7{,}6$\,km/s)~\cite{mohan2024}.

O trabalho realizado na primeira parte deste projeto (POC~I) estabeleceu uma
solução baseline baseada no algoritmo REINFORCE com uma rede MLP de 32
\textit{features}, integrada ao pipeline Hypatia + ns-3. Os resultados
demonstraram que o agente treinado é capaz de reduzir em 47,1\% as
interrupções de conectividade em simulações de 2~horas (mais de um ciclo
orbital completo), com ganho de 4,4\% na latência mediana. Contudo, dois
pontos de melhoria críticos foram identificados:
\begin{enumerate}
  \item O agente realiza 2,65$\times$ mais handovers que a heurística de
        distância mínima, sugerindo que a função de recompensa atual
        ($\beta = 15$) não penaliza adequadamente as trocas desnecessárias.
  \item A MLP processa cada passo de tempo de forma independente, sem
        explorar a periodicidade orbital de $\approx 96$~minutos.
\end{enumerate}

Esta proposta descreve as atividades a serem realizadas na POC~II, que visam
superar essas limitações por meio de três eixos de investigação:
\begin{enumerate}[label=\textbf{\arabic*.}]
  \item \textbf{Reward Shaping Multiobjetivo} --- reformulação da função de
        recompensa para equilibrar qualidade do enlace, disponibilidade de
        conectividade e custo de handover, inspirada em trabalhos recentes de
        handover multiobjetivo em LEO~\cite{song2025}.
  \item \textbf{Memória Temporal via LSTM} --- substituição da MLP por uma
        arquitetura recorrente capaz de capturar dependências temporais nos
        padrões orbitais, abordagem estudada em~\cite{fan2025}.
  \item \textbf{Algoritmo PPO} --- substituição do REINFORCE pelo Proximal
        Policy Optimization (PPO)~\cite{schulman2017}, algoritmo de política
        mais estável e eficiente em termos de amostras.
\end{enumerate}

Adicionalmente, todos os experimentos serão executados com validação estatística
robusta usando múltiplas seeds, seguindo as diretrizes de~\cite{henderson2018}.

% ════════════════════════════════════════════════════════════════════════════
\section{Objetivos}

\subsection{Objetivo Geral}

Estender e aprofundar o sistema de handover inteligente desenvolvido na POC~I,
investigando o impacto isolado e combinado de reward shaping multiobjetivo,
memória temporal (LSTM) e do algoritmo PPO sobre as métricas de desempenho de
rede em uma constelação LEO simulada de 625 satélites.

\subsection{Objetivos Específicos}

\begin{itemize}
  \item \textbf{(H1 --- Reward Shaping)} Avaliar se uma função de recompensa
        multiobjetivo, com pesagem sistemática das componentes de sinal,
        latência, disponibilidade e penalidade de handover, reduz o excesso de
        trocas sem sacrificar significativamente a disponibilidade de conexão.

  \item \textbf{(H2 --- Temporalidade)} Verificar se uma política LSTM melhora
        a capacidade de antecipação orbital em relação à política MLP,
        especialmente antes do primeiro ciclo orbital completo ($\approx
        96$~min).

  \item \textbf{(H3 --- PPO)} Comparar REINFORCE e PPO em termos de
        convergência, estabilidade de treinamento e qualidade da política final
        nas configurações MLP e LSTM.

  \item \textbf{(Escala)} Realizar um experimento exploratório com a constelação
        em escala maior, motivado pelo contexto Starlink com 4.236 satélites
        estudado em~\cite{wu2025}.

  \item \textbf{(Validação)} Executar todos os experimentos com $N=10$ seeds
        independentes e reportar $\mu \pm \sigma$ com intervalos de confiança de
        95\%, seguindo as recomendações de~\cite{henderson2018}.
\end{itemize}

\subsection{Contribuições Esperadas}

\begin{enumerate}
  \item Análise sistemática do trade-off entre continuidade de conexão,
        qualidade de enlace e frequência de handover por meio de reward shaping.
  \item Demonstração empírica do benefício de memória temporal na política de
        handover LEO antes e ao longo do ciclo orbital.
  \item Estudo de ablation comparando arquitetura (MLP vs.\ LSTM) e algoritmo
        (REINFORCE vs.\ PPO) de forma cruzada.
  \item Framework de avaliação estatisticamente robusto e reproduzível para
        políticas de handover em simuladores LEO.
\end{enumerate}

% ════════════════════════════════════════════════════════════════════════════
\section{Referencial Teórico}

\subsection{Handover Multiobjetivo em Redes LEO}

O trabalho de Song et al.~\cite{song2025} formula o handover como um problema
de otimização multiobjetivo, balanceando simultaneamente frequência de handover,
carga dos satélites, throughput e estabilidade do serviço (p.~1--2 do artigo).
Os autores propõem uma nova função de recompensa combinada:
\begin{equation}
  R = w_1 Q_{\text{signal}}
    + w_2 Q_{\text{latency}}
    + w_3 Q_{\text{availability}}
    - w_4 I_{\text{HO}}
    - w_5 \text{Load},
  \label{eq:reward-multi}
\end{equation}
onde cada $w_i$ é um peso a ser ajustado sistematicamente. A
Equação~\eqref{eq:reward-multi} serve de base para o Experimento~1 desta
proposta, transformando o ajuste de $\beta$ em uma investigação científica:
\textit{qual é o trade-off ótimo entre continuidade de conexão, qualidade do
enlace e número de handovers?}

\subsection{Arquitetura LSTM para Previsão e Handover em LEO}

Fan et al.~\cite{fan2025} propõem o uso de LSTM para prever padrões futuros de
tráfego em redes LEO e, com base nessa previsão, tomar decisões proativas de
handover (p.~2--4 do artigo). A arquitetura combina LSTM com atenção e Rainbow
DQN para capturar o estado temporal da rede. Para este projeto, o foco está na
camada recorrente como substituta da MLP:
\begin{equation}
  [s_{t-T+1}, \ldots, s_t] \xrightarrow{\text{LSTM}} h_t \xrightarrow{\pi} a_t,
\end{equation}
onde $T$ é o tamanho da janela temporal. Isso permite ao agente observar a
trajetória orbital, e não apenas o estado instantâneo.

\subsection{Mecanismo de Atenção --- Transformer}

Vaswani et al.~\cite{vaswani2017} introduziram o Transformer, que substitui a
recorrência por mecanismos de self-attention, permitindo modelar relações entre
diferentes posições da sequência sem dependência temporal estrita (p.~2--3 do
artigo). Para o problema de handover LEO, o agente poderia aprender quais
momentos do histórico orbital são mais relevantes para a decisão atual.

\subsection{Proximal Policy Optimization (PPO)}

Schulman et al.~\cite{schulman2017} propõem o PPO, um algoritmo de gradiente de
política que melhora a estabilidade e a eficiência de amostra do REINFORCE
mediante uma função objetivo com razão de probabilidade \textit{clipped} (p.~1--2
do artigo):
\begin{equation}
  L^{\text{CLIP}}(\theta) = \hat{\mathbb{E}}_t\!\left[
    \min\!\left(r_t(\theta)\hat{A}_t,\,
    \operatorname{clip}(r_t(\theta),\,1-\varepsilon,\,1+\varepsilon)\,\hat{A}_t
    \right)
  \right],
  \label{eq:ppo}
\end{equation}
onde $r_t(\theta) = \pi_\theta(a_t|s_t) / \pi_{\theta_{\text{old}}}(a_t|s_t)$.
O PPO permite múltiplas atualizações sobre os dados coletados sem divergência da
política.

\subsection{Escalabilidade: Mega-constelações}

Wu et al.~\cite{wu2025} apresentam o SaTE para a constelação Starlink com 4.236
satélites, usando grafos heterogêneos e redes neurais de grafos para inferência
em tempo real com latência média de 17\,ms (p.~1--3 do artigo). O trabalho
justifica a motivação para escalar os experimentos e contextualiza a dificuldade
computacional associada.

\subsection{Reprodutibilidade em Deep Reinforcement Learning}

Henderson et al.~\cite{henderson2018} demonstram que resultados de Deep RL
apresentam variação significativa devido ao não determinismo e às características
estocásticas do treinamento, recomendando práticas rigorosas de avaliação (p.~2--5
do artigo): múltiplas seeds independentes, testes de significância, e reporte de
$\mu \pm \sigma$ com IC de 95\%.

% ════════════════════════════════════════════════════════════════════════════
\section{Metodologia}

\subsection{Ponto de Partida: Baseline da POC~I}

O ponto de partida é o agente REINFORCE + MLP de 32 features treinado na POC~I,
avaliado no pipeline Hypatia + ns-3 com a constelação de 625 satélites. A função
de recompensa atual é:
\begin{equation}
  r_t = \text{SNR} - \alpha L_{\text{prop}} - \beta I_{\text{HO}} - \delta I_{\text{disconnect}},
  \quad \beta = 15,\; \delta = 50.
  \label{eq:reward-baseline}
\end{equation}
Métricas de referência: número de handovers, interrupções, RTT mediano, UDP
delivery ratio, TCP FCT e throughput.

\subsection{Experimento~1 --- Reward Shaping Multiobjetivo}

Mantendo arquitetura MLP e algoritmo REINFORCE, substitui-se
Eq.~\eqref{eq:reward-baseline} pela função multiobjetivo Eq.~\eqref{eq:reward-multi}.
Varia-se sistematicamente:
\[
  \beta \in \{5,\, 10,\, 15,\, 25,\, 50,\, 75\}.
\]
Analisa-se o trade-off entre handovers e disponibilidade. Os pesos ótimos
definidos aqui serão usados nos experimentos subsequentes.

\subsection{Experimento~2 --- Memória Temporal com LSTM}

Com a recompensa otimizada, substitui-se a MLP por uma LSTM de janela $T$:
\begin{equation}
  \mathbf{S}_t = [s_{t-T+1},\, \ldots,\, s_{t-1},\, s_t] \in \mathbb{R}^{T \times 32}.
\end{equation}
Comparação: REINFORCE + MLP vs.\ REINFORCE + LSTM ($T \in \{10, 20\}$), nos
horizontes de 20\,min e 2\,h.

Pergunta científica:
\begin{quote}
  \textit{A memória temporal permite ao agente explorar a periodicidade orbital
  de $\approx 96$~min e melhorar o desempenho antes do primeiro ciclo orbital
  completo?}
\end{quote}

\subsection{Experimento~3 --- Transformer como Alternativa ao LSTM}

Substitui-se a LSTM por um Transformer Encoder de $L$ camadas sobre a janela
de $T$ passos:
\begin{equation}
  [s_{t-T+1}, \ldots, s_t]
  \xrightarrow{\text{Transformer Encoder}} h_t
  \xrightarrow{\pi} a_t.
\end{equation}
Comparação direta com os resultados do Experimento~2, avaliando desempenho e
custo computacional.

\subsection{Experimento~4 --- Ablation: Arquitetura $\times$ Algoritmo (PPO)}

\begin{table}[H]
  \centering
  \caption{Matriz de ablation.}
  \label{tab:ablation}
  \begin{tabularx}{0.80\textwidth}{lXX}
    \toprule
    \textbf{Configuração}  & \textbf{Arquitetura} & \textbf{Algoritmo} \\
    \midrule
    Baseline (POC~I)       & MLP                  & REINFORCE          \\
    Experimento 4a         & MLP                  & PPO                \\
    Experimento 4b         & LSTM                 & REINFORCE          \\
    Experimento 4c         & LSTM                 & PPO                \\
    \addlinespace
    Referência             & ---                  & Heurística dist.\ mínima \\
    \bottomrule
  \end{tabularx}
\end{table}

\subsection{Experimento~5 --- Validação Estatística}

Todas as configurações da Tabela~\ref{tab:ablation} são executadas com $N=10$
seeds independentes. Reportam-se~\cite{henderson2018} (p.~3--4):
\begin{equation}
  \bar{x} \pm t_{0{,}025,\, N-1} \cdot \frac{s}{\sqrt{N}}.
\end{equation}

\subsection{Experimento~6 (Exploratório) --- Escalabilidade}

Experimento exploratório escalando a constelação para um número maior de
satélites, motivado por Wu et al.~\cite{wu2025} com 4.236 satélites (p.~2
do artigo). O estado do agente mantém $K=4$ candidatos independentemente do
tamanho da constelação, sugerindo boa escalabilidade da política aprendida.

% ════════════════════════════════════════════════════════════════════════════
\section{Hipóteses de Pesquisa}

\begin{description}[leftmargin=2.2cm, style=sameline]
  \item[\textbf{H1 --- Reward Shaping:}]
        Uma função de recompensa multiobjetivo reduz o excesso de handovers sem
        sacrificar significativamente a disponibilidade de conexão.

  \item[\textbf{H2 --- Temporalidade:}]
        Uma política LSTM melhora a capacidade de antecipação em relação à MLP,
        especialmente antes/ao longo do primeiro ciclo orbital completo de
        $\approx 96$~min.

  \item[\textbf{H3 --- Eficiência Algorítmica:}]
        O PPO converge mais rapidamente e produz políticas com menor variância
        do que o REINFORCE, em ambas as arquiteturas (MLP e LSTM).
\end{description}

% ════════════════════════════════════════════════════════════════════════════
\section{Resultados Esperados}

\begin{enumerate}
  \item \textbf{Função de recompensa otimizada}: identificação dos pesos $w_i$
        que minimizam o número de handovers com perda aceitável de disponibilidade.
  \item \textbf{Ganho de antecipação orbital}: demonstração de que a LSTM
        melhora métricas de longo prazo (2\,h) em relação à MLP, especialmente
        antes do primeiro ciclo orbital completo.
  \item \textbf{Estudo de ablation}: tabela comparativa para as quatro
        combinações arquitetura $\times$ algoritmo, com IC de 95\% e análise de
        significância estatística.
  \item \textbf{Resultados estatisticamente robustos}: diferenças reportadas com
        significância estatística, seguindo Henderson et al.~\cite{henderson2018}.
  \item \textbf{Framework reproduzível}: código e scripts disponibilizados para
        extensão futura.
\end{enumerate}

% ════════════════════════════════════════════════════════════════════════════
\section{Cronograma}

\begin{table}[H]
  \centering
  \caption{Cronograma de atividades --- POC~II (10 semanas).}
  \label{tab:cronograma}
  \begin{tabularx}{\textwidth}{cp{4.5cm}p{7.5cm}}
    \toprule
    \textbf{Semana} & \textbf{Atividade Principal} & \textbf{Tarefas Detalhadas} \\
    \midrule
    1
    & Revisão bibliográfica e setup
    & Leitura aprofundada de \cite{song2025,fan2025,schulman2017,henderson2018,wu2025,vaswani2017};
      revisão do código da POC~I; configuração do ambiente de experimentos. \\
    \addlinespace
    2
    & Implementação do Exp.~1 (Reward Shaping)
    & Implementação da nova função de recompensa multiobjetivo (Eq.~\ref{eq:reward-multi});
      definição do grid de pesos $\beta$; primeiros treinamentos. \\
    \addlinespace
    3
    & Avaliação do Exp.~1
    & Execução completa com o grid de pesos; análise das métricas;
      seleção da configuração ótima de recompensa. \\
    \addlinespace
    4
    & Implementação do Exp.~2 (LSTM)
    & Implementação da arquitetura LSTM no Gymnasium;
      janelas $T \in \{10, 20\}$; primeiros treinamentos. \\
    \addlinespace
    5
    & Avaliação do Exp.~2 e início do Exp.~3 (Transformer)
    & Comparação MLP vs.\ LSTM nos horizontes de 20\,min e 2\,h;
      análise da antecipação orbital; início da implementação do Transformer. \\
    \addlinespace
    6
    & Avaliação do Exp.~3 (Transformer)
    & Conclusão da implementação do Transformer;
      comparação LSTM vs.\ Transformer;
      seleção da melhor arquitetura com memória temporal. \\
    \addlinespace
    7
    & Implementação e avaliação do Exp.~4 (PPO)
    & Implementação do PPO; execução do ablation \{MLP, LSTM\} $\times$ \{REINFORCE, PPO\};
      coleta das curvas de convergência. \\
    \addlinespace
    8
    & Validação estatística (Exp.~5) e exploratório (Exp.~6)
    & Execução com $N=10$ seeds; cálculo de $\mu \pm \sigma$ e IC 95\%;
      experimento exploratório de escalabilidade. \\
    \addlinespace
    9
    & Análise consolidada dos resultados
    & Geração de gráficos comparativos, tabelas e CDFs;
      análise do ablation; identificação das melhores configurações. \\
    \addlinespace
    10
    & Escrita do relatório final
    & Redação da monografia (introdução, referencial teórico, metodologia,
      resultados e conclusão); revisão e entrega. \\
    \bottomrule
  \end{tabularx}
\end{table}

% ════════════════════════════════════════════════════════════════════════════
\section{Recursos Necessários}

\begin{itemize}
  \item Ambiente de simulação Hypatia + ns-3 já configurado (herdado da POC~I).
  \item Acesso a GPU para treinamento dos modelos LSTM, Transformer e PPO.
  \item Python $\geq 3.10$: \texttt{PyTorch}, \texttt{Gymnasium},
        \texttt{NumPy}, \texttt{SciPy}, \texttt{Matplotlib}, \texttt{Pandas}.
  \item Dataset de trajetórias TLE da constelação Starlink (herdado da POC~I).
\end{itemize}

% ════════════════════════════════════════════════════════════════════════════
\section*{Referências}
\addcontentsline{toc}{section}{Referências}

\begin{thebibliography}{99}

\bibitem{song2025}
M.~Song, J.~Tian, H.~Zhang, T.~Xu, H.~Hu e T.~Zhou,
``Deep Reinforcement Learning Based Multi-Objective Handover for LEO Satellite Networks,''
em \textit{2025 IEEE 101st Vehicular Technology Conference (VTC2025-Spring)},
IEEE, 2025. DOI: 10.1109/VTC2025-SPRING65109.2025.11174574.
\textit{Citadas p.~1--2.}

\bibitem{fan2025}
D.~Fan, S.~Zhou, J.~Luo, Z.~Yang e M.~Zeng,
``Joint Traffic Prediction and Handover Design for LEO Satellite Networks with LSTM and Attention-Enhanced Rainbow DQN,''
\textit{Electronics}, vol.~14, n.~15, p.~3040, 2025.
DOI: 10.3390/electronics14153040.
\textit{Citadas p.~2--4.}

\bibitem{vaswani2017}
A.~Vaswani, N.~Shazeer, N.~Parmar, J.~Uszkoreit, L.~Jones, A.~N.~Gomez,
L.~Kaiser e I.~Polosukhin,
``Attention Is All You Need,''
em \textit{Advances in Neural Information Processing Systems (NeurIPS)}, 2017.
arXiv:1706.03762.
\textit{Citadas p.~2--3.}

\bibitem{schulman2017}
J.~Schulman, F.~Wolski, P.~Dhariwal, A.~Radford e O.~Klimov,
``Proximal Policy Optimization Algorithms,''
arXiv:1707.06347, OpenAI, 2017.
\textit{Citadas p.~1--2.}

\bibitem{wu2025}
H.~Wu, Y.~Han, M.~Rajpal, Q.~Zhang e J.~Wang,
``SaTE: Low-Latency Traffic Engineering for Satellite Networks,''
em \textit{Proceedings of the ACM SIGCOMM 2025 Conference}, SIGCOMM '25,
pp.~896--916. ACM, 2025. DOI: 10.1145/3718958.3750524.
\textit{Citadas p.~1--3.}

\bibitem{henderson2018}
P.~Henderson, R.~Islam, P.~Bachman, J.~Pineau, D.~Precup e D.~Meger,
``Deep Reinforcement Learning that Matters,''
em \textit{Proceedings of the 32nd AAAI Conference on Artificial Intelligence (AAAI-18)},
2018. DOI: 10.1609/aaai.v32i1.11694.
\textit{Citadas p.~2--5.}

\bibitem{mohan2024}
N.~Mohan, A.~E.~Ferguson, H.~Cech, R.~Bose, P.~R.~Renatin, M.~K.~Marina e J.~Ott,
``A Multifaceted Look at Starlink Performance,''
em \textit{Proceedings of the ACM Web Conference 2024}, WWW '24,
pp.~2723--2734. ACM, 2024.

\end{thebibliography}

\end{document}
